%% appendices/appendixA.tex
%% ============================================================================
%% APPENDIX A: SOURCE CODE
%% ============================================================================
%%
%% PURPOSE: Include complete code that's too long for the main chapters
%%
%% TIPS:
%% - Only include essential code that supports your thesis
%% - Add comments explaining what each section does
%% - Consider linking to a GitHub repository instead of including all code
%% - Format code consistently
%%

\chapter{Source Code}
\label{app:code}


\section{[Module/Component Name]}
\label{sec:app_module1}

% ═══════════════════════════════════════════════════════════════════════════
% EXAMPLE: Complete class or module listing
% ═══════════════════════════════════════════════════════════════════════════

\begin{lstlisting}[language=Python, caption={[Descriptive Name]}, label={lst:app_code1}]
"""
Module: [module_name].py
Purpose: [Brief description of what this code does]
Author: [Your Name]
"""

class ExampleClass:
    """
    A sample class demonstrating the implementation.
    
    Attributes:
        attribute1: Description of first attribute
        attribute2: Description of second attribute
    """
    
    def __init__(self, param1, param2):
        """Initialize the class with given parameters."""
        self.attribute1 = param1
        self.attribute2 = param2
    
    def method_one(self, input_data):
        """
        Process the input data.
        
        Args:
            input_data: The data to process
            
        Returns:
            Processed result
        """
        # Implementation here
        result = self._helper_method(input_data)
        return result
    
    def _helper_method(self, data):
        """Private helper method."""
        # Helper implementation
        return data
\end{lstlisting}


\section{[Another Module/Component]}
\label{sec:app_module2}

\begin{lstlisting}[language=Java, caption={[Descriptive Name]}, label={lst:app_code2}]
/**
 * Class: ExampleService.java
 * Purpose: [Brief description]
 * @author [Your Name]
 */
public class ExampleService {
    
    private final Repository repository;
    
    public ExampleService(Repository repository) {
        this.repository = repository;
    }
    
    /**
     * Main processing method.
     * @param id The entity ID to process
     * @return The processed result
     */
    public Result process(Long id) {
        Entity entity = repository.findById(id);
        // Processing logic here
        return new Result(entity);
    }
}
\end{lstlisting}


\section{Configuration Files}
\label{sec:app_config}

% ═══════════════════════════════════════════════════════════════════════════
% EXAMPLE: Configuration file
% ═══════════════════════════════════════════════════════════════════════════

\begin{lstlisting}[language=Python, caption={Configuration Settings}, label={lst:app_config}]
# config.py - Application configuration

# Database settings
DATABASE = {
    'host': 'localhost',
    'port': 5432,
    'name': 'myapp_db',
    'user': 'app_user'
}

# API settings
API_CONFIG = {
    'base_url': 'https://api.example.com',
    'timeout': 30,
    'retry_count': 3
}

# Feature flags
FEATURES = {
    'enable_caching': True,
    'debug_mode': False,
    'max_items': 100
}
\end{lstlisting}


\section{Database Schema}
\label{sec:app_schema}

% ═══════════════════════════════════════════════════════════════════════════
% EXAMPLE: SQL schema
% ═══════════════════════════════════════════════════════════════════════════

\begin{lstlisting}[language=SQL, caption={Database Schema}, label={lst:app_schema}]
-- Database schema for [Your Application]

CREATE TABLE users (
    id SERIAL PRIMARY KEY,
    username VARCHAR(50) UNIQUE NOT NULL,
    email VARCHAR(100) UNIQUE NOT NULL,
    created_at TIMESTAMP DEFAULT CURRENT_TIMESTAMP
);

CREATE TABLE items (
    id SERIAL PRIMARY KEY,
    user_id INTEGER REFERENCES users(id),
    title VARCHAR(200) NOT NULL,
    description TEXT,
    status VARCHAR(20) DEFAULT 'active',
    created_at TIMESTAMP DEFAULT CURRENT_TIMESTAMP
);

CREATE INDEX idx_items_user_id ON items(user_id);
\end{lstlisting}

% TIP: Add more sections as needed for your code
% TIP: Only include code that's directly relevant to your thesis
